% --- Άσκηση 34322 (34322.tex) ---
\Askhsh{34322}{4}
\begin{ylh}{2}
    \item 3.3. Εξισώσεις 2ου Βαθμού
    \item 4.1. Ανισώσεις 1ου Βαθμού
\end{ylh}
Μια υπολογιστική μηχανή έχει προγραμματιστεί έτσι ώστε, όταν εισάγεται σε αυτήν ένας πραγματικός αριθμός $x$, να δίνει ως εξαγόμενο τον αριθμό $\lambda$ που προκύπτει από τη σχέση: $\lambda = (2x + 5)^{2} - 8x$ (1)
\begin{alist}
\item Αν ο εισαγόμενος αριθμός $x$ είναι ο $- 5$, ποιος είναι ο εξαγόμενος αριθμός $\lambda$;
\monades{4}
\item Αν ο εξαγόμενος αριθμός $\lambda$ είναι ο $20$, ποιος είναι ο εισαγόμενος αριθμός $x$;
\monades{6}
\item
\begin{rlist}
\item
  Να δείξετε ότι η σχέση (1) μπορεί ισοδύναμα να γραφεί στη μορφή:
\begin{equation*}4x^{2} + 12x + (25 - \lambda) = 0.\end{equation*}
\monades{2}
\item
  Να αποδείξετε ότι οποιαδήποτε τιμή και να έχει ο εισαγόμενος αριθμός $x$, ο εξαγόμενος αριθμός $\lambda$ δεν μπορεί να είναι ίσος με $5$.
\monades{6}
\item
  Να προσδιορίσετε τις δυνατές τιμές που μπορεί να έχει ο εξαγόμενος αριθμός $\lambda$.
\monades{7}
\end{rlist}
\end{alist}
